\kafkasection{White Dwarf III:}
\kafkasubsection{What is the thermal energy stored in a white dwarf of temperature $T = \qty{1e7}{K}$ and a mass of $\qty{1}{M_\solar}$?}
The thermal energy is given by:
\begin{equation}
    U = \frac{3}{2}N k_B T  = \frac{3}{2}\frac{\qty{1}{M_\solar}}{\qty{12}{m_p}} k_B T
\end{equation}
\begin{equation}
    U = \frac{3}{2}\frac{\qty{1.989e33}{g}}{12(\qty{1.673e-24}{g})} (\qty{1.381e-16}{erg/K})(\qty{1e7}{K})
\end{equation}
\begin{equation}
    \boxed{U = \qty{2.05e47}{erg}}
\end{equation}
\kafkasubsection{Assuming that the white dwarf radiates like a black body of temperature $T = \qty{1e4}{K}$, estimate its lifetime as a luminous object.}
Let's find the luminosity.
\begin{equation}
    L = 4\pi R^2 \sigma T^4
\end{equation}
Where $T$ is the black body temperature. The age estimate is then how long this luminosity runs through the internal energy as the white dwarf does not produce any more energy. let's now find the radius:
\begin{equation}
    R = \frac{\hbar^2}{Gm_e}\ab(\frac{3Z}{4\pi A m_p})^{5/3} M^{-1/3}
\end{equation}
\begin{equation}
    R = \frac{(\qty{6.626e-27}{erg\ s})^2}{4\pi^2 (\qty{6.673e-8}{dyne\ cm^2 /g^2})(\qty{9.110e-28}{g})}\ab(\frac{3(6)}{4\pi(12)(\qty{1.673e-24}{g})})^{5/3} \ab(\qty{1.989e33}{g})^{-1/3}
\end{equation}
\begin{equation}
    R = \qty{1.79e7}{cm}
\end{equation}
Okay now we're just supposed to throw this radius out and give the more accurate answer derived from the figure in the textbook.
\begin{equation}
    R = \qty{0.6e9}{cm}
\end{equation}
The luminosity is then:
\begin{equation*}
    L = 4\pi (\qty{0.6e9}{cm})^2 (\qty{5.670e-5}{erg/cm^2/K^4/s}) (\qty{1e4}{K})^4
\end{equation*}
\begin{equation}
    L = \qty{2.57e30}{erg/s}
\end{equation}
So now the luminous time is:
\begin{equation}
    \frac{U}{L} = \frac{\qty{2.05e47}{erg}}{\qty{2.57e30}{erg/s}} = \boxed{\qty{7.98e16}{s}}
\end{equation}
This is in the billions of years.